\documentclass[10pt]{article}
\usepackage[margin=0.7in]{geometry}
\usepackage{booktabs}
\usepackage{url}
\title{Time without an accepted alive state:\\a bounded offline Lightpool remeasurement}
\author{Codex (AI-assisted draft)}
\date{22 November 2026 --- unreviewed; not approved for publication}
\begin{document}
\maketitle
\begin{abstract}
Dead-declaration recovery omits suspicions that resolve before death. We add
lifetime-bound accounting of time without an accepted alive membership state,
then remeasure a frozen set of 24 saved virtual attempts without executing the
protocol. The two source baselines remain separate. The six attempts containing
the exchange-origin correction are null; the older dense workload has
suspect-only burden. These measurements establish an outcome capability and
bounded descriptive observations, not candidate benefit or ordinary-network
fidelity.
\end{abstract}

\section{Question and outcome}
How much time do fixed healthy observer--target pairs spend without an accepted
alive state in the retained baselines? SWIM distinguishes suspicion from final
failure dissemination~\cite{swim}; Lifeguard changes suspicion and refutation
timing~\cite{lifeguard}. These motivate observing the interval before a dead
declaration, but neither paper supplies the outcome contract used here.
We make no historical numerical replication or novelty claim.

An episode starts at the first departure from alive, continuing across suspect,
dead, left and absent membership states until accepted alive in the same two
process lifetimes. Removal means absence, not recovery. We retain those four
components separately; time without alive is their sum. It is an operational
membership outcome, not an exact duration of detector mistakes. Packet receipt,
incarnation changes and confirmations do not constitute recovery or restart an
episode. The primary descriptive quantity is healthy-pair time without alive
divided by eligible pair-time, with completion and censor counts alongside it.

\section{Frozen evidence and method}
\texttt{inputs.json} fixes exactly six attempts from the partial-partition record
(two cases, indices 1--3) and 18 virtual companions from the ordinary-witness
record (six cases, indices 1--3). It retains original identities, toolchains,
statuses, exits and all raw payload hashes. No ordinary process, simulation,
seed expansion or replacement attempt runs. The evidence snapshot is pinned by
\texttt{evidence-base}; the measurement implementation is pinned by
\texttt{reducer-base} (\texttt{5dbbca8}), built with Go 1.26.3. Original specs
continue to identify their original producer, not the new reducer.

The partial-partition producer is \texttt{54ddfd4}, including the exchange-origin
correction. Its four A nodes and target T lose both directions for 16.384 s;
27 O-to-T pairs are predeclared healthy. The older producer/tool is
\texttt{62c746a}; its core/wire match \texttt{4b45bcd}, before the correction.
It has 24 A nodes with A-to-T impairment and seven healthy O-to-T pairs. Dense
sensitivity cases vary only the named latency or modeled service duration.
The corresponding healthy controls preserve their respective pair sets.

At W=15 s, immediately before boundary events, all-members-alive warm-up fixes
process lifetimes. Elapsed time covers [15 s,120 s); events at both endpoints
are processed in sequence/change order. Eligible exposure includes pauses and
all beliefs, ending permanently at endpoint crash, restart, shutdown or leave.
Healthy pairs must have no configured direct impairment in either direction or
an overlapping endpoint pause. Pre-window non-alive beliefs invalidate warm-up;
carry-in diagnostic time is clipped and excluded from duration quantiles.
Incomplete prefixes end at their last valid event and remain incomplete.

The reducer streams the trace without invoking protocol code. State-time and
segment conservation, zero-time episodes, removal, carry-in, lifetime ordering,
population selection and censoring have independent synthetic regression tests.
Two separate scratch mutations (admit every healthy pair; reset onset on a
confirmation) each fail its intended test. The recipe validates saved reports,
compares all old summary fields after normalizing only the output format,
checks raw inputs before and after reduction, and reproduces all retained
outputs from the committed record. Summary v2 adds the duration family while
saved input v1 and trace v1 remain unchanged.

\section{Results}
\begin{table}[ht]\centering\footnotesize
\setlength{\tabcolsep}{3pt}
\begin{tabular}{llrrrrr}\toprule
Case & Index & Pair-s & Without alive (ns) & Completed & p95 (ns) & Max (ns) \\ \midrule
PP (new) & 1 & 2835 & 0 & 0 & -- & -- \\
PP (new) & 2 & 2835 & 0 & 0 & -- & -- \\
PP (new) & 3 & 2835 & 0 & 0 & -- & -- \\
PP healthy (new) & 1 & 2835 & 0 & 0 & -- & -- \\
PP healthy (new) & 2 & 2835 & 0 & 0 & -- & -- \\
PP healthy (new) & 3 & 2835 & 0 & 0 & -- & -- \\
Dense (old) & 1 & 735 & 12893499711 & 28 & 655637735 & 667846924 \\
Dense (old) & 2 & 735 & 6583778400 & 27 & 410940862 & 544555961 \\
Dense (old) & 3 & 735 & 9149864678 & 21 & 930121281 & 972587432 \\
Dense healthy (old) & 1 & 735 & 0 & 0 & -- & -- \\
Dense healthy (old) & 2 & 735 & 0 & 0 & -- & -- \\
Dense healthy (old) & 3 & 735 & 0 & 0 & -- & -- \\
Latency 100us (old) & 1 & 735 & 8759888660 & 28 & 555657762 & 637226129 \\
Latency 100us (old) & 2 & 735 & 7402252627 & 28 & 483877971 & 539858255 \\
Latency 100us (old) & 3 & 735 & 9148964678 & 21 & 930121281 & 972587432 \\
Latency 10ms (old) & 1 & 735 & 11309396624 & 28 & 703889275 & 719332316 \\
Latency 10ms (old) & 2 & 735 & 28507056002 & 13 & 4099736802 & 4099736802 \\
Latency 10ms (old) & 3 & 735 & 9123456976 & 21 & 676248645 & 716740375 \\
Service 100us (old) & 1 & 735 & 12764280877 & 28 & 655537735 & 667546924 \\
Service 100us (old) & 2 & 735 & 7769384511 & 25 & 520692714 & 527196228 \\
Service 100us (old) & 3 & 735 & 14521135751 & 27 & 1087191724 & 1100279989 \\
Service 1ms (old) & 1 & 735 & 10575901343 & 26 & 653648927 & 716999110 \\
Service 1ms (old) & 2 & 735 & 7809163172 & 32 & 400000000 & 496309350 \\
Service 1ms (old) & 3 & 735 & 14523496119 & 28 & 761538564 & 796635571 \\
\bottomrule\end{tabular}
\caption{Fixed healthy-pair outcomes, kept separate by source baseline and case. All without-alive time is suspect time. All censor counts are zero; -- denotes an empty completion sample, not a zero quantile.}
\end{table}

All 24 original attempts are complete and valid; there are no exclusions,
carry-ins, lifecycle interruptions or censored healthy-pair episodes. Every
healthy non-alive interval is suspect-only: dead, left and absent components
are each zero in every attempt. The six corrected partial-partition attempts
have zero burden over 2,835 healthy pair-seconds each. The three older healthy
controls also have zero burden over 735 pair-seconds each. Empty duration samples
remain null rather than zero-duration recoveries. Positive older dense cells
have no healthy dead declarations, so dead-only recovery omits their entire
healthy non-alive burden. Table~1 preserves per-attempt denominators and
conditional duration summaries; full p50/p95/max, every pair and ordered
segments are retained in the derived JSON and \texttt{cells.tsv}.

\section{Costs, limitations and reproduction}
This is an offline outcome remeasurement on two distinct virtual baselines.
Older dense positives cannot establish headroom on the corrected baseline;
the corrected null does not authorize a larger search. Three indices per case
support no inferential or tail claim. Observer pairs share protocol histories
and are not independent samples. No survival estimator is fitted. Quantiles
condition on completion and are not censor-adjusted estimates. The modeled
network/service assumptions and the ordinary/virtual mechanism mismatch remain
unresolved; this work neither calibrates ordinary execution nor admits a
candidate campaign.

The record stores losslessly compressed derived JSON, deterministic logs,
checksums, table inputs and this unreviewed paper; original raw evidence is
referenced rather than duplicated. No new resource study was performed.
Runtime and storage for offline reduction are implementation costs, not protocol
traffic or benefit measurements. The implementation keeps per-pair state and
output episodes/segments, never the whole trace.

From a checkout containing the record, run
\path{experiments/2026-11-21-non-alive-belief-duration/run.sh}
with a new output directory.
It restores originals with the pinned evidence utility, reduces one at a time,
inventories all attempts even after a reduction failure, and returns nonzero
for unexpected failures or reproduction differences. \texttt{render.sh NEW\_DIRECTORY}
downloads a checksum-pinned Tectonic 0.15.0 and verifies its fixed TeX bundle
identity before rendering. The rendering toolchain and exact commands are
recorded in \texttt{paper-toolchain.txt}. The operator owns final wording and
any external release.

\begin{thebibliography}{2}
\bibitem{swim} A. Das, I. Gupta, and A. Motivala.
SWIM: Scalable Weakly-consistent Infection-style Process Group Membership
Protocol. DSN, 2002.
\url{https://www.cs.cornell.edu/projects/Quicksilver/public_pdfs/SWIM.pdf}.
\bibitem{lifeguard} A. Dadgar, J. Phillips, and J. Currey.
Lifeguard: Local Health Awareness for More Accurate Failure Detection.
arXiv:1707.00788v2, 2018. \url{https://arxiv.org/abs/1707.00788v2}.
\end{thebibliography}
\end{document}
